\documentclass[aps,prl,twocolumn,showpacs,superscriptaddress]{revtex4-2}
\usepackage{amsmath,amssymb,booktabs,xcolor,hyperref,amsthm}
\newtheorem{theorem}{Theorem}

\begin{document}

\title{Appendix KA1: OHC Resonance Mode Structure and Shell Capacity\\
Full derivation of $2n^2$ capacity from Hopf fibration geometry}

\author{Michel Robert Cabri\'e}
\affiliation{Independent Researcher, Victoria, Australia}
\email{ZeroFreeParameters@gmail.com}
\date{20 March 2026}

\begin{abstract}
This appendix derives the OHC Bessel resonance mode structure and 
the electron shell capacity $2n^2$ from first principles, supporting 
Letter~78 (Atomic Structure as Hopfion Topology). The $n$-th OHC 
resonance level has exactly $n^2$ independent spatial modes, each 
supporting two Hopf orientations ($H=\pm1$), giving total capacity 
$2n^2$. The two BCT void types generate the orbital symmetries: 
$r_{\rm oct}$ gives 1 isotropic (s) mode per shell; $r_{\rm tet}$ 
gives 3 directional (p) modes per shell.
\end{abstract}

\maketitle

\section{Purpose}

This appendix provides the technical derivation supporting Letter~78.
The key results are: (1) shell capacity $2n^2$ from Bessel mode 
counting; (2) orbital symmetry from void geometry; (3) the Hopf 
charge constraint $|H|\leq n$ per mode.

\section{OHC Bessel Mode Derivation}

The OHC interior condensate $\Psi_{\rm int}$ satisfies the 
Gross-Pitaevskii equation within each Planck sphere of radius 
$R = \ell_P/2$ (Letter~27). Linearised modes are spherical Bessel 
resonances indexed by $(n, \ell, m)$ where $n$ is the principal 
quantum number, $\ell$ the orbital angular momentum, and $m$ the 
magnetic quantum number.

For given $n$, the allowed pairs satisfy $0\leq\ell\leq n-1$ and 
$-\ell\leq m\leq+\ell$. The total spatial mode count:
\begin{equation}
    \sum_{\ell=0}^{n-1}(2\ell+1) = 1+3+5+\cdots+(2n-1) = n^2
\end{equation}
This is the sum of the first $n$ odd integers. Each spatial mode 
supports two Hopf orientations ($H=+1$ spin up, $H=-1$ spin down), 
giving total shell capacity $2n^2$. \qed

\section{Mode Capacity and Hopf Charge}

The constraint $|H_{\rm total}|\leq n$ per mode follows from the 
Josephson boundary condition at coupling $\alpha_0 = r_{\rm oct}
\times r_{\rm tet}/\pi = 0.0074081$. The Josephson phase across 
the $n$-th Bessel mode is $\Delta\varphi_n = 2\pi n\alpha_0$. 
The maximum coherent Hopf excitation sustainable within this phase 
budget is $|H|\leq n$.

\begin{center}
\begin{tabular}{cclrr}
\toprule
$n$ & $\ell$ values & Spatial modes & $\times2$ spin & Noble gas \\
\midrule
1 & $\ell=0$ & $n^2=1$ & $2n^2=2$ & He ($Z=2$) \\
2 & $\ell=0,1$ & $n^2=4$ & $2n^2=8$ & Ne ($Z=10$) \\
3 & $\ell=0,1,2$ & $n^2=9$ & $2n^2=18$ & Ar ($Z=18$) \\
4 & $\ell=0,1,2,3$ & $n^2=16$ & $2n^2=32$ & Kr ($Z=36$) \\
\bottomrule
\end{tabular}
\end{center}

\section{Orbital Symmetry from Void Geometry}

The two BCT void types generate the orbital symmetries:
\begin{itemize}
\item $r_{\rm oct} = (\sqrt{2}-1)/2$: 1 isotropic centre 
$\to$ $\ell=0$ (s-orbital), 1 mode per shell
\item $r_{\rm tet} = (\sqrt{6}-2)/4$: 3 axial orientations 
$\to$ $\ell=1$ (p-orbitals), 3 modes per shell
\item Higher Hopf sectors: $\ell=2$ (d, 5 modes), $\ell=3$ (f, 7 modes)
\end{itemize}

Total per shell at $n=4$: $1+3+5+7=16$ spatial modes $\times 2 = 
32 = 2\times4^2$ \checkmark

The octet rule (Letter~78) follows directly: the $n=2$ shell has 
$1~({\rm oct}) + 3~({\rm tet}) = 4$ spatial modes $\times 2 = 8$ electrons.

\begin{thebibliography}{99}
\bibitem{AppJH} M.~R.~Cabri\'e, \textit{Appendix JH: The OHC},
Zenodo (2026), doi:10.5281/zenodo.18975018
\bibitem{L27} M.~R.~Cabri\'e, \textit{BCT Letter 27}, Zenodo (2026).
\bibitem{L78} M.~R.~Cabri\'e, \textit{BCT Letter 78}, Zenodo (2026).
\end{thebibliography}

\end{document}
